\section{Data Processing}

The data layer begins with passenger-flow organization rather than timetable recurrence. This order is essential because the timetable is not an independent object; it is a downstream consequence of the chosen service pattern, and that service pattern is only justified when the burden concentration along the line has first been made visible. Therefore, the data section should first answer where the overlap zone is likely to appear and only then prepare the variables used in the later operation and timetable models.

The first preprocessing step is to aggregate station-level or OD-level passenger information into section burdens that can support service-pattern selection. Depending on the attachment form, this may require converting boarding-alighting counts, OD matrices, or time-sliced demand records into representative section flows over the study horizon. The purpose is not merely to produce one plot, but to reveal whether the central corridor and the terminal sections carry sufficiently different burdens to justify differentiated service.

The second step is to construct operation-planning inputs. Once the overlap zone is identified, the data must be converted into variables that the service-pattern model can actually use, such as required section capacity, feasible headway range, rolling-stock availability, and station dwell assumptions. In a complete paper, these quantities are often summarized in a compact planning table so that the reader can see how the later optimization problem is parameterized.

The third step is to prepare timetable-recurrence and audit inputs. Tracking intervals, station dwell bounds, turnaround assumptions, and scenario demand settings should all be stabilized here before the solution section begins. Doing so keeps the later timetable equations readable and prevents the validation section from introducing seemingly new assumptions without warning. It also lets the reader understand why some operation plans may look attractive in the cost-service comparison but fail under timetable feasibility checks.

By the end of the data section, the reader should be able to answer three questions: where the strong overlap demand lies, which planning parameters govern the large-small-route comparison, and which operating assumptions will later be audited in the timetable stage. If these points are visible, the transition from data to model is no longer abrupt and the timetable route becomes much easier to trust.
